\section{Experiments}
\label{sec:experiments}

The experiments are designed to answer four questions. First, does uncertainty-aware querying maintain task performance while reducing human interaction? Second, is the timing and content of queries important compared with random or fixed query policies? Third, how does the confidence threshold $\tau$ control the trade-off between autonomy and interaction? Fourth, does selective querying reduce human workload compared with continuous supervision or action-level intervention?

\subsection{Domains}

\placeholderfig{Experimental domains used to evaluate the proposed framework, including simulated waste sorting and tomato harvesting with real-robot validation.}{2.0in}{fig:domains}

\paragraph{Waste sorting.}
The simulated waste sorting domain consists of a UR5 robot sorting objects into category-specific bins. Each object belongs to one of four categories: general waste, paper, can, or plastic. Categories are not fully known in advance and may be ambiguous due to occlusion or perception uncertainty. The robot must detect, pick, and place objects into the correct bins.

\paragraph{Tomato harvesting.}
The tomato harvesting domain contains tomatoes attached to stems. Each tomato may be ripe, unripe, or rotten. Ripe tomatoes should be harvested, rotten tomatoes should be discarded, and unripe tomatoes should be left untouched. The real-robot platform consists of a mobile base, a UR5e arm, an RGB-D camera, and a LiDAR sensor. This domain evaluates whether the framework can connect symbolic uncertainty, real perception, human feedback, and robot execution.

\subsection{Baselines and Ablations}

We compare the proposed method against the following strategies:
\begin{itemize}[leftmargin=*]
    \item \textbf{No Query}: the robot never asks for human feedback.
    \item \textbf{Always Query}: the robot queries at every uncertain decision point, approximating full human supervision.
    \item \textbf{Random Query}: the robot queries randomly with the same average query budget as the proposed method.
    \item \textbf{Failure-Triggered Query}: the robot asks only after reaching a failure, dead end, or invalid action.
    \item \textbf{Ours without Frontier}: the robot uses confidence thresholding but selects queries without action-relevant frontier structure.
    \item \textbf{Ours}: the full uncertainty-aware frontier query framework.
\end{itemize}

\subsection{Metrics}

We report task success rate, cumulative reward, average number of queries, query ratio, task completion time, and failure type. For user studies, we additionally report subjective workload using NASA-TLX and interaction preference ratings. For real-robot evaluation, we report scenario-level success, number of queries, intervention timing, and representative failure cases.

\subsection{System Evaluation}

\begin{table}[t]
\centering
\caption{System evaluation across waste sorting and tomato harvesting domains.}
\label{tab:system_eval}
\begin{tabular}{llcccc}
\toprule
Domain & Method & Success (\%) & Reward & Queries & Time (s) \\
\midrule
\multirow{6}{*}{Waste}
& No Query & -- & -- & -- & -- \\
& Always Query & -- & -- & -- & -- \\
& Random Query & -- & -- & -- & -- \\
& Failure-Triggered & -- & -- & -- & -- \\
& Ours w/o Frontier & -- & -- & -- & -- \\
& Ours & -- & -- & -- & -- \\
\midrule
\multirow{6}{*}{Tomato}
& No Query & -- & -- & -- & -- \\
& Always Query & -- & -- & -- & -- \\
& Random Query & -- & -- & -- & -- \\
& Failure-Triggered & -- & -- & -- & -- \\
& Ours w/o Frontier & -- & -- & -- & -- \\
& Ours & -- & -- & -- & -- \\
\bottomrule
\end{tabular}
\end{table}

The results compare task performance and interaction cost across query strategies. Always Query is expected to achieve high success at the cost of frequent interaction, whereas No Query may fail when ambiguity remains unresolved. Random Query often spends queries on non-critical states. The proposed method aims to approach the task success of Always Query with fewer human queries by targeting uncertainty that is relevant to planning.

\subsection{Threshold Analysis}

\placeholderfig{Threshold sensitivity for $\tau$ across waste sorting and tomato harvesting domains.}{2.0in}{fig:tau}

The confidence threshold $\tau$ controls the autonomy-interaction trade-off. Low $\tau$ values make the robot more willing to commit uncertain beliefs, reducing queries but increasing failure risk. High $\tau$ values make the robot ask more often, improving confidence but increasing workload. We analyze success rate, reward, and query count across $\tau$ values in both domains because $\tau$ is the key parameter governing selective autonomy.

\subsection{Real-Robot Tomato Evaluation}

The real-robot study includes repeated trials across multiple tomato scenarios. Each scenario specifies the initial arrangement, tomato states, visual ambiguities, and expected task outcome. The evaluation reports whether each query prevents a wrong harvest, wrong discard, or missed target. A qualitative case study compares the same scene under No Query and Ours, showing how a human answer changes the belief and the subsequent plan.

\subsection{User Study}
\label{sec:user_study}

The user study evaluates whether system-initiated state queries reduce workload compared with continuous supervision and action-level assistance. Participants observe or interact with the robot through the interface while the robot executes task scenarios containing perceptual ambiguity. We compare at least three interaction conditions: \textbf{Always Query}, where the user is repeatedly asked to verify uncertain information; \textbf{Failure-Triggered Query}, where the user is asked only after the robot reaches an invalid or failed state; and \textbf{Ours}, where the system asks targeted state-level questions when entropy-based confidence is low.

We measure subjective workload using NASA-TLX, along with the number of user responses, response time, perceived interruption, trust, and preference. The main hypothesis is that the proposed method reduces workload because the user is not required to continuously monitor the robot or make action-level planning decisions. Instead, the user answers targeted questions about symbolic state facts, while the robot remains responsible for belief update and replanning.

\begin{table}[t]
\centering
\caption{User study results comparing interaction conditions.}
\label{tab:user_study}
\begin{tabular}{lcccc}
\toprule
Condition & NASA-TLX & Responses & Response Time & Preference \\
\midrule
Always Query & -- & -- & -- & -- \\
Failure-Triggered & -- & -- & -- & -- \\
Action-Level Query & -- & -- & -- & -- \\
Ours & -- & -- & -- & -- \\
\bottomrule
\end{tabular}
\end{table}
